\section{Arquitectura y representación de memoria}
\label{sec:arquitectura}

\subsection{Ejecutables y módulos}

La compilación produce los cuatro ejecutables del enunciado
(\code{creator}, \code{producer}, \code{consumer} y \code{finisher}) y un
quinto, \code{verify}, que cruza los CSV de una corrida
(Sección~\ref{sec:pruebas}). Cada productor y consumidor es un proceso
independiente que solo conoce el nombre de la instancia. El código común se
divide en los cinco módulos que pide el enunciado y unos pocos de apoyo:

\begin{itemize}
  \item \textbf{Memoria} (\code{memory.c}): crea, valida, mapea y elimina
  el segmento.
  \item \textbf{Buffer} (\code{buffer.c}): arreglo circular puro, sin
  sincronización; sus funciones exigen el mutex ya tomado.
  \item \textbf{Sincronización} (\code{sync.c}): la única capa que toca el
  mutex y las variables de condición; implementa \code{put}, \code{get} y
  la pausa interrumpible.
  \item \textbf{Participantes} (\code{participants.c}): conexión, registro,
  retiro y máquina de estados del cierre.
  \item \textbf{Métricas} (\code{metrics.c}): agregados de la instancia y
  resumen del finalizador.
  \item \textbf{Apoyo}: análisis de argumentos (\code{cli.c}), CSV
  (\code{logger.c}), carga útil verificable (\code{payload.c}), generador
  \emph{xorshift32}~\cite{Marsaglia2003} (\code{rng.c}) y reloj
  monotónico (\code{clock.c}).
\end{itemize}

\subsection{Un segmento por instancia}

Cada instancia es \emph{un único} objeto de memoria compartida POSIX con el
nombre que da el usuario (por ejemplo \code{/soa\_p2\_g1}, visible en Linux
como \code{/dev/shm/soa\_p2\_g1})~\cite{POSIX2024}. No hay semáforos con
nombre ni archivos auxiliares: el mutex y las variables de condición viven
\emph{dentro} del segmento, así que dos instancias con nombres distintos no
comparten ningún recurso. El nombre debe ser \code{/} seguido de 1 a 30
caracteres de \code{[A-Za-z0-9\_-]}.

El segmento contiene un único \code{struct instance}, con el arreglo de
mensajes como miembro flexible al final. La Fig.~\ref{fig:segmento}
muestra su disposición con los desplazamientos medidos en x86-64 con
glibc. La cabecera ocupa 12\,528~bytes y cada mensaje 64, así que el
segmento mide $12\,528 + 64\cdot\mathit{capacidad}$~bytes: 13\,040 con
capacidad 8 y unos 78~KB con 1024.

\begin{figure}[!t]
\centering
\begin{tikzpicture}[
    font=\scriptsize, node distance=0pt,
    campo/.style={draw, text width=5.4cm, align=center, inner xsep=2pt,
                  inner ysep=2.5pt, fill=black!4},
    off/.style={anchor=east, text=black!70},
    prot/.style={anchor=west, text=black!70},
    m/.style={draw, minimum height=0.42cm, inner sep=0pt, font=\tiny,
              fill=black!12, anchor=north west}
  ]
  \node[campo] (c0) {\code{magic} $\cdot$ \code{version} $\cdot$ \code{capacity} $\cdot$ \code{state}};
  \node[campo, below=of c0] (c1) {\code{mutex} (40 B) $+$ \code{not\_full}, \code{not\_empty},
    \code{closing}, \code{all\_left} (48 B c/u)};
  \node[campo, below=of c1] (c2) {\code{head}, \code{tail}, \code{count},
    \code{enqueue\_seq}, \code{dequeue\_seq}, \code{active\_producers},
    \code{active\_consumers}};
  \node[campo, below=of c2] (c3) {\code{refs}: conexiones $+$ bit de puerta cerrada};
  \node[campo, below=of c3] (c4) {\code{producers}: 64 casillas de registro de 64 B};
  \node[campo, below=of c4] (c5) {\code{consumers}: 64 casillas de registro de 64 B};
  \node[campo, below=of c5] (c6) {\code{metrics}: marcas de tiempo, rechazos e
    histograma de latencia (496 cubetas)};
  \node[campo, below=of c6, fill=black!12] (c7) {\code{messages[0..capacidad$-$1]}, 64 B c/u};
  \foreach \i/\o/\p in {0/0/{A, I, I, M}, 1/16/primitivas, 2/248/M, 3/288/A,
                        4/296/M, 5/4\,400/M, 6/8\,504/M, 7/12\,528/M} {
    \node[off]  at (c\i.west) {\o\ };
    \node[prot] at (c\i.east) {\ \p};
  }
  \node[font=\scriptsize\bfseries, anchor=south east] at (c0.north west) {offset};
  \node[font=\scriptsize\bfseries, anchor=south west] at (c0.north east) {protec.};

  % Detalle de un mensaje (no a escala)
  \coordinate (m0) at ($(c7.south west)+(0,-0.45)$);
  \node[m, minimum width=0.72cm] (m1) at (m0) {\code{prod\_id}};
  \node[m, minimum width=0.55cm] (m2) at (m1.north east) {\code{resv}};
  \node[m, minimum width=0.85cm] (m3) at (m2.north east) {\code{sequence}};
  \node[m, minimum width=0.85cm] (m4) at (m3.north east) {\code{enq\_seq}};
  \node[m, minimum width=0.85cm] (m5) at (m4.north east) {\code{enq\_ns}};
  \path let \p1=(m5.east), \p2=(c7.east) in
    node[m, minimum width=\x2-\x1] (m6) at (m5.north east) {\code{payload} (32 B)};
  \draw[densely dotted] (c7.south west) -- (m1.north west);
  \draw[densely dotted] (c7.south east) -- (m6.north east);
  \node[font=\tiny, anchor=north west, align=left, text width=5.6cm, inner xsep=0pt]
    at ($(m1.south west)+(0,-0.05)$)
    {\raggedright payload: \code{sequence} y \code{producer\_id} (copias),
     \code{checksum} \mbox{FNV-1a} y \code{fill[16]}};
\end{tikzpicture}
\caption{Disposición de \code{struct instance} en el segmento (x86-64, glibc
2.39). Protección: A = atómico, I = inmutable tras publicar, M = mutex de
la instancia. Abajo, el detalle de un mensaje de 64~B (no a escala).}
\label{fig:segmento}
\end{figure}

\subsection{Datos compartidos y datos privados}

Todo lo que dos procesos necesitan ver igual está en el segmento: el estado
de la instancia, las primitivas, los índices y contadores del buffer, las
tablas de registro de por vida y los agregados para el resumen final.

Cada proceso guarda en su memoria privada la \emph{dirección} donde mapeó
el segmento (en general distinta entre procesos), el estado de su
generador pseudoaleatorio sembrado una vez con \code{--seed}, su
\code{FILE*} de CSV, sus contadores de espera y adquisición (que copia a su
casilla de registro al retirarse) y el mensaje que construye o procesa.
Así, el procesamiento simulado y la escritura del CSV no necesitan el
mutex.

Ningún campo compartido contiene punteros: el buffer usa los índices
\code{head} y \code{tail} módulo la capacidad sobre \code{messages[]}, y
los mensajes son valores de tamaño fijo que se copian con
\code{buffer\_push} y \code{buffer\_pop}. La carga útil es verificable sin
estado extra: $\mathit{fill}[i] = (\mathit{producer\_id} +
\mathit{sequence} + i) \bmod 256$, y un \emph{checksum} FNV-1a de
32~bits~\cite{FNV} sobre la serialización campo por campo de
\code{sequence}, \code{producer\_id} y \code{fill}. El consumidor la
recalcula al extraer y anota el resultado en su CSV.
Los atómicos del segmento (\code{magic} y \code{refs}) solo son válidos
entre procesos si no usan un candado interno, lo que se exige en tiempo de
compilación con \code{ATOMIC\_INT\_LOCK\_FREE == 2}.

\subsection{Creación exclusiva y publicación}

El \code{creator} sigue un orden fijo:
\begin{enumerate}
  \item valida el nombre y la capacidad;
  \item llama a \code{shm\_open} con \code{O\_CREAT|O\_EXCL};
  \item ajusta el tamaño con \code{ftruncate}, reintentando ante
  \code{EINTR};
  \item mapea el segmento con \code{mmap};
  \item inicializa las primitivas con \code{PTHREAD\_PROCESS\_SHARED} y las
  variables de condición con \code{CLOCK\_MONOTONIC};
  \item escribe \code{version}, \code{capacity} y \code{state = OPEN};
  \item por último, escribe \code{magic} con \emph{release}.
\end{enumerate}

\code{O\_EXCL} hace la creación atómica: entre varios \code{creator}
simultáneos gana uno, y los demás reciben \code{EEXIST} y salen
\emph{sin tocar} el segmento ajeno. Si falla un paso posterior, el
\code{creator} deshace solo lo que creó: destruye las primitivas
inicializadas, desmapea y hace \code{shm\_unlink} del nombre que acaba de
reservar.

Quien se conecta lee \code{magic} con \emph{acquire} y rechaza un
segmento más chico que la cabecera, un \code{magic} distinto de la
constante (\code{creator} sin terminar o segmento ajeno) y una versión,
capacidad o tamaño incompatibles. El par \emph{release}/\emph{acquire} garantiza que quien observa
\code{magic} observa también todo lo escrito antes. Como el kernel entrega
el segmento en ceros y la constante es distinta de cero, una instancia a
medio crear nunca parece lista.

Un mutex no puede proteger esta publicación: quien llega no sabe todavía
si el mutex está inicializado.

\subsection{Propiedad de los recursos}

Cada recurso de una instancia tiene un único responsable de crearlo y
otro de liberarlo, y esa asignación es la que evita fugas y dobles
destrucciones. La Tabla~\ref{tab:propiedad} la resume; el detalle de qué
campo protege cada mecanismo está en \code{include/memory.h}.

\begin{table}[H]
\centering
\caption{Propiedad de los recursos de una instancia}
\label{tab:propiedad}
\footnotesize
\begin{tabular}{@{}p{2.3cm}p{2.5cm}p{2.9cm}@{}}
\toprule
Recurso & Crea & Libera \\
\midrule
Nombre del segmento & \code{creator} (\code{O\_EXCL}) & \code{finisher} (\code{shm\_unlink}) \\
Objeto de memoria & \code{creator} (\code{ftruncate}) & el kernel, al último \code{munmap} \\
Mapeo de un proceso & ese proceso (\code{mmap}) & ese proceso (\code{munmap}) \\
Mutex y condvars & \code{creator}, antes de publicar & \code{finisher}, con $\mathit{refs}=1$ \\
Referencia de conexión & \code{instance\_attach} & \code{instance\_detach} \\
Casilla de registro & \code{participant\_\allowbreak register} & nunca: queda \code{LEFT} \\
CSV de cada proceso & ese proceso & ese proceso \\
\bottomrule
\end{tabular}
\end{table}

Casi todos los recursos los crea el \code{creator} o el propio proceso que
los usa, y los libera el \code{finisher} o ese mismo proceso. Hay dos
excepciones. El objeto de memoria no lo libera nadie del programa: el
\code{finisher} solo elimina el \emph{nombre}, y el kernel libera el
objeto cuando el último proceso lo desmapea. Esa propiedad vuelve seguro
el rechazo de una conexión tardía (Sección~\ref{sec:protocolo}). La
casilla de registro, por su parte, no se libera nunca: queda marcada como
retirada para que su id no pueda volver a usarse. Las primitivas se
destruyen solo cuando el \code{finisher} es el único conectado
($\mathit{refs}=1$).
